\documentclass[uplatex,a4paper,11pt]{jsarticle}

\usepackage[dvipdfmx]{graphicx}
\usepackage{amsmath, amssymb}
\usepackage{geometry}
\usepackage{hyperref}
\usepackage{here}
\usepackage{bookmark}

\geometry{margin=25mm}
\renewcommand{\baselinestretch}{1.2}

\title{Homework 4a}
\author{学籍番号: 1270374 \\ 名前: 松本 奈生}
\date{\today}

\begin{document}
\maketitle

%% ============================================================
\section*{Question 1}
%% ============================================================
\subsection*{(a) $\mathtt{a^*b^*}$}

最初の5要素：
$\varepsilon,\ \mathtt{a},\ \mathtt{b},\ \mathtt{aa},\ \mathtt{ab}$

要素でない最初の5文字列：
$\mathtt{ba},\ \mathtt{baa},\ \mathtt{bab},\ \mathtt{bba},\ \mathtt{bbb}$

%% ------------------------------------------------------------
\subsection*{(b) $\mathtt{a(ba)^*b}$}

最初の5要素：
$\mathtt{ab},\ \mathtt{abab},\ \mathtt{ababab},\ \mathtt{abababab},\ \mathtt{ababababab}$

要素でない最初の5文字列：
$\varepsilon,\ \mathtt{a},\ \mathtt{b},\ \mathtt{aa},\ \mathtt{ba}$

%% ------------------------------------------------------------
\subsection*{(e) $\Sigma^*\mathtt{a}\Sigma^*\mathtt{b}\Sigma^*\mathtt{a}\Sigma^*$}

少なくとも $\mathtt{a}, \mathtt{b}, \mathtt{a}$ をこの順に含む文字列の集合。

最初の5要素：
$\mathtt{aba},\ \mathtt{aaba},\ \mathtt{abaa},\ \mathtt{abab},\ \mathtt{abba}$

要素でない最初の5文字列：
$\varepsilon,\ \mathtt{a},\ \mathtt{b},\ \mathtt{aa},\ \mathtt{ab}$

%% ------------------------------------------------------------
\subsection*{(g) $(\varepsilon \cup \mathtt{a})\mathtt{b}$}

最初の5要素（全要素）：
$\mathtt{b},\ \mathtt{ab}$

要素でない最初の5文字列：
$\varepsilon,\ \mathtt{a},\ \mathtt{aa},\ \mathtt{ba},\ \mathtt{bb}$

%% ============================================================
\section*{Question 2}
%% ============================================================

\subsection*{1. $M_2$ を入力 $\mathtt{000000}$ に適用したときの configuration の系列}

$M_2$ は入力文字列が $\mathtt{0}$ の個数が偶数かどうかを判定する機械である。
以下では状態 $q_i$ とテープ内容・ヘッド位置を configuration として表す。
（ヘッド位置の直前に状態を記載する形式を用いる。）

\begin{align*}
&q_1\mathtt{000000} \\
\vdash\; &\sqcup q_2\mathtt{00000} \\
\vdash\; &\sqcup x\,q_3\mathtt{0000} \\
\vdash\; &\sqcup x\mathtt{0}\,q_4\mathtt{000} \\
\vdash\; &\sqcup x\mathtt{0}x\,q_3\mathtt{00} \\
\vdash\; &\sqcup x\mathtt{0}x\mathtt{0}\,q_4\mathtt{0} \\
\vdash\; &\sqcup x\mathtt{0}x\mathtt{0}x\,q_3\sqcup \\
\vdash\; &\sqcup x\mathtt{0}x\mathtt{0}\,q_5 x\sqcup \\
\vdash\; &\sqcup x\mathtt{0}x\,q_5\mathtt{0}x\sqcup \\
\vdash\; &\sqcup x\mathtt{0}\,q_5 x\mathtt{0}x\sqcup \\
\vdash\; &\sqcup xq_5 \mathtt{0}\,x\mathtt{0}x\sqcup \\
\vdash\; &\sqcup q_5 x\mathtt{0}\,x\mathtt{0}x\sqcup \\
\vdash\; &\sqcup x\,q_2\mathtt{0}x\mathtt{0}x\sqcup \\
\vdash\; &\sqcup xx\,q_3 x\mathtt{0}x\sqcup \\
\vdash\; &\sqcup xxx\,q_3\mathtt{0}x\sqcup \\
\vdash\; &\sqcup xxx\mathtt{0}\,q_4 x\sqcup \\
\vdash\; &\sqcup xxx\mathtt{0}x\,q_4\sqcup \\
\vdash\; &\sqcup xxx\mathtt{0}x\,q_{reject}
\end{align*}

%% ------------------------------------------------------------
\subsection*{2. $M_1$ を入力 $\mathtt{10\#11}$ に適用したときの configuration の系列}

$M_1$ は $\{w\#w \mid w \in \{0,1\}^*\}$ を決定する機械である。
入力 $\mathtt{10\#11}$ は $w=\mathtt{10}$，$w'=\mathtt{11}$ であり $w \neq w'$ なので拒否される。

\begin{align*}
&q_1\,\mathtt{10\#11} \\
\vdash\; &x\,q_3\,\mathtt{0\#11} \\
\vdash\; &x\mathtt{0}\,q_3\,\mathtt{\#11} \\
\vdash\; &x\mathtt{0\#}\,q_5\,\mathtt{11} \\
\vdash\; &x\mathtt{0}\,q_6\,\mathtt{\#x1} \\
\vdash\; &x\,q_6\,\mathtt{0\#x1} \\
\vdash\; &q_6\,x\mathtt{0\#x1} \\
\vdash\; &x\,q_{reject}\,\mathtt{0\#x1}
\end{align*}


%% ============================================================
\section*{Question 3}
%% ============================================================

言語 $L = \{w \mid w \text{ は } \mathtt{1} \text{ の2倍の個数の } \mathtt{0} \text{ を持つ}\}$
を決定するチューリング機械 $M$ の実装レベルの記述を以下に示す。

\noindent
$M =$入力文字列 $w$ に対して：

\begin{enumerate}
  \item テープ上を左から右へ走査し，まだ処理されていない $1$ を1つ見つける．
  そのような $1$ が存在する場合は，それを記号 $X$ に書き換える．
  そのような $1$ が存在しない場合は，手順4へ進む．

  \item テープの左端に戻り，左から右へ走査して，まだ処理されていない $0$ を1つ見つける．
  そのような $0$ が存在する場合は，それを記号 $Y$ に書き換える．
  存在しない場合は，入力を拒否せよ．

  \item 再びテープの左端に戻り，左から右へ走査して，まだ処理されていない $0$ を1つ見つける．
  そのような $0$ が存在する場合は，それを記号 $Y$ に書き換える．
  存在しない場合は，入力を拒否せよ．
  その後，手順1へ戻る．

  \item テープ上を左から右へ走査し，まだ処理されていない $0$ または $1$ が存在するかを確認する．
  もし存在する場合は，入力を拒否せよ．
  存在しない場合は，入力を受理せよ．
\end{enumerate}

%% ============================================================
\section*{Question 4}
%% ============================================================

\noindent\textbf{命題：}
$L_D = \{\langle M \rangle \mid M \text{ は } \langle M \rangle \text{ を受理しないチューリング機械}\}$
はチューリング認識不可能である。

\noindent
言語
\[
L_D = \{ \langle M \rangle \mid M \text{ は } \langle M \rangle \text{ を受理しない} \}
\]
が Turing 認識可能であると仮定する．
すなわち，$L_D$ を認識する Turing 機械 $D$ が存在すると仮定する．

このとき，$\langle D \rangle$ が $L_D$ に属するかどうかで場合分けを行う．

\begin{enumerate}
  \item $\langle D \rangle \in L_D$ の場合：

  $L_D$ の定義より，$D$ は $\langle D \rangle$ を受理しない．

  一方で，$D$ は $L_D$ を認識するので，
  $L_D$ の要素はすべて受理しなければならない．
  よって $\langle D \rangle \in L_D$ ならば，
  $D$ は $\langle D \rangle$ を受理するはずである．

  これは矛盾である．

  \item $\langle D \rangle \notin L_D$ の場合：

  $L_D$ の定義より，$D$ は $\langle D \rangle$ を受理する．

  しかし，$D$ は $L_D$ を認識するので，
  $L_D$ に属さない入力は受理してはならない．
  よって $\langle D \rangle \notin L_D$ ならば，
  $D$ は $\langle D \rangle$ を受理しないはずである．

  これも矛盾である．
\end{enumerate}

以上より，いずれの場合にも矛盾が生じるため，
$L_D$ を認識する Turing 機械 $D$ は存在しない．

したがって，$L_D$ は Turing 認識不可能である．

\end{document}
